% !TeX root = main.tex
\section{Results: Downstroke-Ratio Sweep Near Trim}
\subsection{Sweep Setup}
We sweep the downstroke ratio $\delta$ over $[0.30,0.70]$ in increments of $0.02$ while keeping the baseline operating condition fixed.
For each $\delta$, we compute cycle-averaged thrust $\bar{F}_x$, lift $\bar{F}_z$, input power $\bar{P}_{in}$, and propulsive efficiency
\begin{equation}
\bar{\eta} = \frac{\bar{F}_x U}{\bar{P}_{in}}.
\end{equation}

\subsection{Representative Numerical Results}
Table~\ref{tab:delta_points} reports representative values for three downstroke ratios.
The symmetric baseline $\delta=0.5$ yields a trim-like near-zero mean thrust condition.
Reducing $\delta$ increases both thrust and lift at the cost of increased power, while increasing $\delta$ can increase thrust but tends to reduce lift.

\begin{table}[t]
\centering
\caption{Representative duty-cycle sweep results (baseline: $f=3.8~\mathrm{Hz}$, $U=7.0~\mathrm{m/s}$, body pitch $13^\circ$, separation off, prescribed twist on).}
\label{tab:delta_points}
\begin{tabular}{@{}cccccc@{}}
\toprule
$\delta$ & $\bar{F}_x$ (N) & $\bar{F}_z$ (N) & $\bar{P}_{in}$ (W) & $\bar{\eta}$ & Notes \\
\midrule
0.30 & 0.655 & 10.126 & 55.63 & 0.082 & higher thrust/lift, higher power \\
0.50 & 0.095 & 9.209 & 42.56 & 0.016 & symmetric baseline \\
0.70 & 0.588 & 8.736 & 50.50 & 0.081 & higher thrust, lower lift \\
\bottomrule
\end{tabular}
\end{table}

\subsection{Control-Relevant Interpretation}
Duty-cycle modulation induces a strong and nonlinear mapping from $\delta$ to cycle-averaged force and power.
This provides a control-relevant degree of freedom: the vehicle can select $\delta$ to emphasize thrust, lift, or power depending on task requirements.
In particular, for forward-speed control, $\delta$ can serve as a direct input to modulate the mean streamwise force while remaining within a constrained kinematic family.
